\documentclass[11pt]{article}
\usepackage[a4paper,margin=24mm]{geometry}
\usepackage{amsmath,amssymb,booktabs,array,xcolor}
\usepackage[colorlinks=true,linkcolor=blue!45!black,urlcolor=blue!45!black]{hyperref}
\setlength{\parindent}{0pt}
\setlength{\parskip}{6pt}
\setlength{\emergencystretch}{2em}
\newcommand{\Q}{\mathbb Q}
\newcommand{\R}{\mathbb R}
\newcommand{\Z}{\mathbb Z}
\newcommand{\dd}{\,d}
\newcommand{\Bcal}{\mathcal B}
\newcommand{\Hfive}[1]{H_{#1}^{(5)}}
\newcommand{\source}[1]{{\small\color{blue!45!black}Lai: #1.}\par}
\newcommand{\pagetitle}[2]{\section*{#1}\source{#2}}
\pagestyle{plain}
\begin{document}
\begin{center}
{\LARGE Parallel constructions in Lai's sections 1--7}\\[8pt]
{\large One functional, one bilinear form, several representations}\\[6pt]
Explanatory draft \enspace$\cdot$\enspace 14 November 2026
\end{center}

This draft accompanies sections 1--7 of Chun Wing Lai's
\href{https://ai.vixra.org/pdf/2609.0080v1.pdf}{\emph{On Fauzan's Proof that
$\zeta(5)$ is Irrational}}, version 1, dated 29 September 2026.
All source equation numbers below refer to that version. The purpose is to
explain the relationships among the constructions introduced there.
The Stieltjes-transform and bilinear-form language supplies an organizing
interpretation of those formulas.

\pagetitle{The organizing picture}{sections 1--3; definitions in sections 4--7}
The apparent proliferation of objects becomes easier to follow when we ask
two questions: \emph{Which expressions describe the same functional?}
And \emph{which matrices describe the same pairing?}
Write
\[
 \xi=\zeta(5),\qquad K=40n,\qquad N=3n,\qquad h=K-N=37n,
 \qquad n\in\Z_{\geq1},
\]
and put
\[
 D_m(t)=\prod_{j=1}^{m}(t+j^2),\qquad
 T(t)=\frac{D_K(t)}{D_N(t)},\qquad W(t)=D_N(t)^5.
\]
The main relationships are
\[
 \boxed{\mu_X(R)=\tau_X\!\left(x^5R(-x^2)\right)},\qquad
 \boxed{\Bcal_X(p,q)=\mu_X\!\left(\frac{Wpq}{T}\right)}.
\]
Here $p,q$ range over polynomials of degree less than $h$. The first identity
relates two presentations of a functional. The second packages the matrix
entries into one symmetric bilinear form. The parameter $X$ is separate from
the rational-function variables $t$ and $x$.

\begin{center}
\renewcommand{\arraystretch}{1.22}
\begin{tabular}{@{}p{.38\textwidth}p{.55\textwidth}@{}}
\toprule
Parallel descriptions & What the relationship explains \\
\midrule
Bernoulli formulas / positive moments & Algebraic values become integrals at $X=\xi$. \\
$\mu_X$ / $\tau_X$ & A weighted substitution changes square poles into integer poles. \\
$D_N^6/D_K$ / $W/T$ & Cancellation exposes the surviving poles. \\
Hankel / Gram / affine matrix & Index structure, positivity, and dependence on $X$ describe one matrix. \\
Continuous / discrete pairings & Integration proves positivity; pole evaluations prove exact degree. \\
$\Delta_K$ / $F_K$ & A positive scalar changes normalization while preserving roots and degree. \\
\bottomrule
\end{tabular}
\end{center}

Sections 1--3 give the setting and the target: integer polynomials of degree
$O(n)$ whose values at $\xi$ are nonzero and decay like $e^{-cn^2}$.
The constructions below explain how the degree and nonvanishing fit together.
Integer normalization with a sufficiently small value requires estimates
beyond section 7; those estimates are outside this draft.

\newpage
\pagetitle{One functional: algebraic values and positive moments}
{section 4, equations (11)--(26)}
Let $\Hfive{j}=\sum_{v=1}^{j}v^{-5}$, with $\Hfive{0}=0$, and use
Bernoulli numbers with $B_1=-1/2$. A rational function with only simple poles
at negative integer squares has a unique finite decomposition
\[
 R(t)=P(t)+\sum_{j\geq1}\frac{c_j}{t+j^2}.
\]
It therefore suffices to specify a linear functional on monomials and simple
pole terms. The definitions are
\begin{align*}
 \mu_X(t^e)&=\frac{(-1)^e B_{2e+2}(2e+3)(2e+4)(2e+5)}{24},\\
 \mu_X\!\left(\frac1{t+j^2}\right)
 &=j^4\bigl(X-\Hfive{j}\bigr)-\frac14+\frac1{2j}.
\end{align*}
The polynomial part is independent of $X$; only the pole part varies.
For example, $\mu_X(1)=5/12$ and $\mu_X((t+1)^{-1})=X-3/4$.

At $X=\xi$, these two clauses agree with a single positive integral:
\[
 \mu_\xi(R)=\int_0^\infty R(y^2)w(y)\dd y,\qquad
 w(y)=\frac{(2\pi)^4y^5}{12}\sum_{\ell\geq1}\ell^4e^{-2\pi\ell y}>0.
\]
The integral converges for all rational functions in the stated domain:
they have no poles on $[0,\infty)$, the weight stays bounded near zero,
and its exponential decay controls polynomial growth at infinity.
For a monomial, integration term by term gives
\[
 \int_0^\infty y^{2e}w(y)\dd y
 =\frac{(2e+5)!}{12(2\pi)^{2e+2}}\zeta(2e+2).
\]
Euler's formula for the even zeta value turns this into the Bernoulli formula.
Thus the signs and factorials encode moments of a positive weight.

The pole values have a complementary description. Define
\[
 \Phi(a)=\int_0^\infty\frac{w(y)}{a^2+y^2}\dd y
 =a^4\zeta(5,a)-\frac1{2a}-\frac14,\qquad a>0.
\]
At $a=j$, the identity
$\zeta(5,j)=\xi-\Hfive{j}+j^{-5}$ gives exactly
$\Phi(j)=\mu_\xi((t+j^2)^{-1})$.
In the variable $z=a^2$, $\Phi(\sqrt z)$ is the Stieltjes transform of
$\dd\nu(t)=w(\sqrt t)\dd t/(2\sqrt t)$.

This same transform also contains the polynomial moments. A finite geometric
expansion, for any integer $m\geq1$, gives
\[
 \Phi(a)=\sum_{e=0}^{m-1}\frac{(-1)^e\mu_\xi(t^e)}{a^{2e+2}}
 +\frac{(-1)^m}{a^{2m}}\int_0^\infty
       \frac{y^{2m}w(y)}{a^2+y^2}\dd y.
\]
The remainder has absolute value at most $\mu_\xi(t^m)a^{-2m-2}$.
The two definitions of $\mu_\xi$ are therefore linked as integer samples
and asymptotic coefficients of one transform. This is a finite remainder
identity, not an assertion that the resulting infinite series converges.

\newpage
\pagetitle{The Bernoulli functional and the weighted pullback}
{section 7, equations (49)--(59), compared with section 4}
Define $L(x^k)=B_k$ and
\[
 \tau(P)=\frac{L(P''')}{24},\qquad
 \tau_X\!\left(\frac1{x-r}\right)=\Hfive{d(r)}-X,
 \qquad
 d(r)=\begin{cases}r,&r\geq0,\\-r-1,&r<0.\end{cases}
\]
Extend linearly to polynomial parts and simple integer poles. The third
derivative is essential: it supplies three consecutive factors and lowers
the Bernoulli index by three.

To see the Bernoulli summation structure explicitly, let $\partial=d/dx$.
For a polynomial $P$,
\[
 \tau(P)=\frac1{24}\sum_{k\geq0}\frac{B_k}{k!}P^{(k+3)}(0).
\]
This is a finite sum. Its formal symbol is
$u^4/(24(e^u-1))$, obtained from the Bernoulli generating function
$u/(e^u-1)$ by multiplying by $u^3/24$. The familiar Bernoulli summation
operator thus explains the polynomial part of $\tau$ without adding a new
functional.

Now define the linear map
\[
 (\mathcal T R)(x)=x^5R(-x^2).
\]
The claim is $\mu_X=\tau_X\circ\mathcal T$. It can be checked on the two
parts of the partial-fraction basis.

For $R(t)=t^e$, the transformed expression is $(-1)^ex^{2e+5}$, so
\[
 \tau_X\bigl((-1)^ex^{2e+5}\bigr)
 =\frac{(-1)^e(2e+5)(2e+4)(2e+3)B_{2e+2}}{24}
 =\mu_X(t^e).
\]
For a pole, polynomial division and partial fractions give
\[
 \frac{x^5}{j^2-x^2}
 =-x^3-j^2x-\frac{j^4}{2}
       \left(\frac1{x-j}+\frac1{x+j}\right).
\]
Since $\tau(-x^3)=-1/4$, $\tau(-j^2x)=0$, $d(j)=j$ and $d(-j)=j-1$,
\begin{align*}
 \tau_X\!\left(\frac{x^5}{j^2-x^2}\right)
 &=-\frac14-\frac{j^4}{2}
       \left(\Hfive{j}+\Hfive{j-1}-2X\right)\\
 &=j^4\bigl(X-\Hfive{j}\bigr)-\frac14+\frac1{2j}.
\end{align*}
The last step uses $\Hfive{j}-\Hfive{j-1}=j^{-5}$.
Linearity proves the identity for every admissible $R$.

This is a weighted pullback: both $t=-x^2$ and the factor $x^5$ matter.
The pole $t=-j^2$ becomes the pair $x=j,-j$. The apparent asymmetry of
$d(r)$ supplies the harmonic increment that produces $1/(2j)$.
The $t$-description exposes positivity on $t=y^2\geq0$; the
$x$-description exposes integer pole locations.

\newpage
\pagetitle{One modifier and one bilinear form}
{section 5, equations (27), (31)--(38)}
The modifier in the matrix entries has two equal expressions:
\[
 A(t):=\frac{D_N(t)^6}{D_K(t)}
       =\frac{D_N(t)^5}{T(t)}=\frac{W(t)}{T(t)},\qquad
 T(t)=\prod_{k=N+1}^K(t+k^2).
\]
The first expression displays the original construction. The second makes
cancellation explicit: poles at $-1^2,\ldots,-N^2$ disappear, leaving exactly
$K-N=h$ simple poles. The zeros at the canceled locations have order five.

For $E=\Q[t]_{<h}$ define
\[
 \Bcal_X(p,q)=\mu_X\bigl(A(t)p(t)q(t)\bigr),\qquad p,q\in E.
\]
The rational function inside $\mu_X$ has only the allowed simple poles.
This form is symmetric and bilinear. In the basis $1,t,\ldots,t^{h-1}$,
its matrix is
\[
 G_K(X)=\bigl[\Bcal_X(t^i,t^j)\bigr]_{0\leq i,j<h}
       =\bigl[\mu_X(A(t)t^{i+j})\bigr]_{0\leq i,j<h}.
\]
Three descriptions of this matrix now follow from three ways of reading
the same formula.

\textbf{Hankel structure.} The entry depends on $i+j$, so entries are constant
along each anti-diagonal. This is a statement about the monomial coordinates.

\textbf{Gram structure.} At $X=\xi$, extend the form to real polynomials.
The integral representation gives
\[
 \Bcal_\xi(p,q)=\int_0^\infty p(y^2)q(y^2)A(y^2)w(y)\dd y.
\]
Thus $G_K(\xi)$ is the Gram matrix of $1,y^2,\ldots,y^{2h-2}$ for the
positive measure $A(y^2)w(y)\dd y$. A nonzero polynomial has only finitely
many zeros, so $\Bcal_\xi(p,p)>0$ for $p\ne0$. Consequently
$G_K(\xi)$ is positive definite and $\det G_K(\xi)>0$.

\textbf{Affine dependence on $X$.} Write
\[
 A(t)t^{i+j}=P_{ij}(t)+\sum_{k=N+1}^K\frac{c_k^{ij}}{t+k^2},
 \qquad
 c_k^{ij}=\frac{W(-k^2)(-k^2)^{i+j}}{T'(-k^2)}.
\]
The polynomial term contributes no $X$, and every pole contributes linearly.
It follows that $G_K(X)=G_0+XG_1$. Such a family is called an affine matrix
pencil. Here $G_0=G_K(0)$, while $(G_1)_{ij}=\sum_k c_k^{ij}k^4$.
The polynomial part must be retained: it is zero when $i+j+6N-K<0$,
and otherwise has degree $i+j+6N-K$.

The pullback gives a further expression for this same form:
\[
 \Bcal_X(p,q)=\tau_X\!\left(x^5 A(-x^2)p(-x^2)q(-x^2)\right).
\]
No additional determinant is introduced by changing functional notation.

\newpage
\pagetitle{Continuous and discrete evaluations of the matrix}
{sections 5--6, equations (36), (38)--(48)}
The Gram description uses integration over $y>0$. A parallel description
uses evaluation at the remaining poles. It describes the coefficient of
$X$, rather than the whole positive form.

For $R=P+\sum_j c_j/(t+j^2)$ put
\[
 \Lambda(R)=\sum_j c_jj^4.
\]
Then the definitions immediately give
\[
 \mu_X(R)=\mu_\xi(R)+(X-\xi)\Lambda(R),\qquad
 G_K(X)=G_K(\xi)+(X-\xi)G_1.
\]
The single functional $\Lambda$ records all dependence on the parameter.
It also equals the coefficient of $t^{-3}$ in the Laurent expansion of
$R(t)$ at infinity, since
$1/(t+j^2)=t^{-1}-j^2t^{-2}+j^4t^{-3}+\cdots$.

Set $a_k=-k^2$ and
\[
 d_k=\frac{k^4W(a_k)}{T'(a_k)},\qquad N<k\leq K.
\]
The coefficient form is a finite evaluation pairing:
\[
 \Lambda(Apq)=\sum_{k=N+1}^K d_kp(a_k)q(a_k).
\]
Compare this with
$\Bcal_\xi(p,q)=\int p(y^2)q(y^2)A(y^2)w(y)\dd y$.
Both pairings multiply two polynomial values and then combine them with
weights. Their supports and signs differ: the integral has positive weights
on the positive axis, whereas the finite pairing has nonzero, generally
signed weights at negative squares. It is not a quadrature formula for
$\Bcal_\xi$.

Let $V_{ik}=a_k^i$, with $0\leq i<h$ and $N<k\leq K$. Then
\[
 G_1=V\operatorname{diag}(d_k)V^{\mathsf T}.
\]
There are exactly $h$ distinct nodes, so $V$ is a square invertible
Vandermonde matrix. Every $d_k$ is nonzero. Therefore $G_1$ is invertible and
\[
 [X^h]\det G_K(X)=\det G_1\ne0,\qquad \deg\det G_K(X)=h.
\]
The determinant can be evaluated explicitly. Since
\[
 \prod_{k=N+1}^K T'(a_k)
 =(-1)^{h(h-1)/2}\prod_{N<k<\ell\leq K}(a_\ell-a_k)^2,
\]
the squared Vandermonde factors cancel, leaving
\[
 \boxed{\det G_1=(-1)^{h(h-1)/2}
       \prod_{k=N+1}^K k^4D_N(-k^2)^5\in\Z\setminus\{0\}.}
\]
The two evaluation pictures establish different facts: continuous positivity
gives a positive value at $\xi$; discrete nondegeneracy gives exact degree.
The equality between matrix size and number of remaining poles is what
makes the latter argument work so directly.

\newpage
\pagetitle{One determinant, two normalizations}
{section 5, equations (28)--(30); section 3, Lemma 3.1}
Define
\[
 \Delta_K(X)=\det G_K(X),\qquad F_K(X)=S_K\Delta_K(X),
\]
where
\[
 S_K=\frac{(K!)^{2h}4^{h-1}}
 {(N!)^{12h}\displaystyle\prod_{i=1}^{h-1}((2i)!)^2}>0.
\]
These are scalar multiples of one polynomial. Since $S_K$ is independent
of $X$ and nonzero, $F_K$ and $\Delta_K$ have the same roots, with the
same multiplicities, and the same degree $h$. Since $S_K>0$, both have
positive values at $\xi$.

Normalization does change the size of the coefficients and the value at
$\xi$. In particular, the integer leading coefficient found above does not
imply that every coefficient of $\Delta_K$ is integral. Nor does the displayed
definition of $S_K$ alone prove $F_K\in\Z[X]$. Those are arithmetic questions
about the full polynomial, not consequences of its degree or positivity.

There is a related coordinate fact. If the basis of $E$ is changed by an
invertible rational matrix $C$, with new basis vectors given by its columns,
the matrix becomes $C^{\mathsf T}G_K(X)C$. Hence
\[
 \det\bigl(C^{\mathsf T}G_K(X)C\bigr)
 =\det(C)^2\Delta_K(X).
\]
This explains why the bilinear form is the intrinsic object, while a
determinant also records the scale of the chosen coordinates. The displayed
$S_K$ is a specified scalar normalization; no identification of it with a
particular change of basis is needed here.

\textbf{How this meets the opening criterion.}
The criterion in section 3 asks for $Q_n\in\Z[X]$ with degree at most $37n$
and $0<|Q_n(\xi)|<e^{-cn^2}$ for some $c>0$. If $\xi=a/b$, then
\[
 0<\bigl|b^{37n}Q_n(a/b)\bigr|
 <\exp\bigl(37n\log b-cn^2\bigr),
\]
which would eventually put a nonzero integer strictly between zero and one
in absolute value. The contradiction uses degree, integrality, nonvanishing,
and decay together.

By the end of the constructions discussed here, the form supplies
nonvanishing, the pole evaluations supply exact degree, and the pullback
connects the two functional languages. A further scalar multiple can always
clear finitely many rational denominators, but its size matters: multiplying
by a large number can destroy a small-value estimate. Controlling this cost
together with decay is the remaining task, beyond the scope of sections 1--7.

\textbf{Reading boundary.}
The draft stops with the identities and constructions above. The continuous
and discrete pairings are related parts of one affine family, rather than
interchangeable positive measures. The identities organize the construction;
they do not substitute for the later quantitative estimates.

\vfill
{\small\textbf{Source.} Chun Wing Lai, \emph{On Fauzan's Proof that
$\zeta(5)$ is Irrational}, v1, 29 September 2026, sections 1--7,
\href{https://ai.vixra.org/pdf/2609.0080v1.pdf}{ai.vixra.org/pdf/2609.0080v1.pdf}.
The mathematical construction is attributed there to Aabir Fauzan.
The transform, coefficient-functional, and coordinate explanations in this
draft are derived reorganizations of the displayed definitions.}
\end{document}
